\subsection{Change of multiplicative subset}

Let A be a ring, S a multiplicative subset of A and M an A-module. If T is a multiplicative subset of A containing S, it follows from Proposition 5 of no. 2 that there exists a unique \(S^{-1}A\) -linear mapping \(i_{M}^{T,S}:S^{-1}M\to T^{-1}M\) such that \(i_{M}^{T}=i_{M}^{T,S}\circ i_{M}^{S}\) ; the mapping \(i_{M}^{T,S}\) maps the element m/s of \(S^{-1}M\) to the element m/s of \(T^{-1}M\) . It is easily verified that \(i_{M}^{T,S}=i_{A}^{T,S}\otimes l_{M}\) . If U is a third multiplicative subset of A such that \(T\subset U\) , then \(i_{M}^{U,S}=i_{M}^{U,T}\circ i_{M}^{T,S}\) ; moreover, if u: \(M\to N\) is an A-module homomorphism, the diagram

\[
\begin{array}{c} \mathrm{S} ^ {- 1} \mathrm{M} \xrightarrow {\mathrm{S} - 1 _ {u}} \mathrm{S} ^ {- 1} \mathrm{N} \\ i _ {\mathrm{M}} ^ {\mathrm{T}, \mathrm{S}} \Biggl \downarrow \qquad \qquad \qquad \qquad \qquad \qquad \qquad \Biggl \downarrow i _ {\mathrm{N}} ^ {\mathrm{T}, \mathrm{S}} \\ \mathrm{T} ^ {- 1} \mathrm{M} \xrightarrow [ \mathrm{T} - 1 _ {u} ]{} \mathrm{T} ^ {- 1} \mathrm{N} \end{array}
\]

is commutative.

PROPOSITION 7. Let A be a ring and S, T two multiplicative subsets of A. Set \(\mathbf{T}' = i(\mathbf{T}_{\mathbf{A}}^{\mathbf{s}})\) .

\begin{enumerate}
\def\labelenumi{(\roman{enumi})}
\tightlist
\item
  There exists a unique isomorphism \(j\) from the ring (ST) \(^{-1}\) A onto the ring \(\mathbf{T}^{\prime -1}(\mathbf{S}^{-1}\mathbf{A})\) such that the diagram

\[
\begin{array}{c} \mathrm{A} \xrightarrow {i _ {\mathrm{A}} ^ {\mathrm{S}}} \mathrm{S} ^ {- 1} \mathrm{A} \\ i _ {\mathrm{M}} ^ {\mathrm{ST}} \Big \downarrow \qquad \qquad \qquad \qquad \qquad \qquad \qquad \Big \downarrow i _ {\mathrm{S} - 1 _ {\mathrm{A}}} ^ {\mathrm{T} ^ {\prime}} \\ (\mathrm{ST}) ^ {- 1} \mathrm{A} \xrightarrow [ j ]{} \mathrm{T} ^ {\prime - 1} (\mathrm{S} ^ {- 1} \mathrm{A}) \end{array}
\]

is commutative.

\item
  Let M be an A-module. There exists an \((\mathrm{ST})^{-1}\mathrm{A}\) -isomorphism k from the \((\mathrm{ST})^{-1}\mathrm{A}\) -module \((\mathrm{ST})^{-1}\mathrm{M}\) onto the \(\mathrm{T}^{\prime-1}(\mathrm{S}^{-1}\mathrm{A})\) -module \(\mathrm{T}^{\prime-1}(\mathrm{S}^{-1}\mathrm{M})\) such that the diagram
\end{enumerate}

\[
\begin{array}{c} \mathrm{M} \qquad \xrightarrow {i _ {\mathrm{M}} ^ {\mathrm{S}}} \mathrm{S} ^ {- 1} \mathrm{M} \\ i _ {\mathrm{A}} ^ {\mathrm{ST}} \Big \downarrow \qquad \qquad \qquad \qquad \qquad \Big \downarrow i _ {\mathrm{S}} ^ {\mathrm{T} ^ {\prime} - 1 _ {\mathrm{M}}} \\ (\mathrm{ST}) ^ {- 1} \mathrm{M} \xrightarrow [ k ]{} \mathrm{T} ^ {\prime - 1} (\mathrm{S} ^ {- 1} \mathrm{M}) \end{array}
\]

is commutative.

\begin{enumerate}
\def\labelenumi{(\roman{enumi})}
\tightlist
\item
  We use the definition of \((\mathrm{ST})^{-1}\mathrm{A}\) as the solution of a universal mapping problem. Let B be a ring and f a homomorphism from A to B such that \(f(\mathrm{ST})\) consists of invertible elements. As \(f(\mathrm{S})\) consequently consists of invertible elements, there exists a unique homomorphism \(f': S^{-1}A \to B\) such that \(f = f' \circ i_{A}^{S}\) (no. 1, Proposition 1). For all \(t \in T\) , \(f'(i_{\mathbf{A}}^{\mathrm{S}}(t)) = f(t)\) is invertible in B by hypothesis, hence \(f'(\mathrm{T}')\) consists of invertible elements; then there exists, by no. 1, Proposition 1, a unique homomorphism \(f''\) from \(T'^{-1}(S^{-1}A)\) to B such that \(f' = f'' \circ i_{S^{-1}A}^{T'}\) , whence \(f = f'' \circ u\) , setting \(u = i_{S^{-1}A}^{T'} \circ i_{A}^{S}\) .

Moreover, if \(f_1'' \colon \mathrm{T}^{-1}(\mathrm{S}^{-1}\mathrm{A}) \to \mathrm{B}\) is a second homomorphism such that \(f_1'' \circ u = f\) , then \((f_1'' \circ i_{\mathrm{S}^{-1}\mathrm{A}}^{\mathrm{T}'} \circ i_{\mathrm{A}}^{\mathrm{S}} = (f'' \circ i_{\mathrm{S}^{-1}\mathrm{A}}^{\mathrm{T}'} \circ i_{\mathrm{A}}^{\mathrm{S}} \text{, whence } f_1'' \circ i_{\mathrm{S}^{-1}\mathrm{A}}^{\mathrm{T}'} = f'' \circ i_{\mathrm{S}^{-1}\mathrm{A}}^{\mathrm{T}'} \text{ and consequently } f_1'' = f''\) .

As the images under \(u\) of the elements of ST in \(\mathbf{T}^{-1}(\mathbf{S}^{-1}\mathbf{A})\) are invertible, the ordered pair \((\mathbf{T}^{-1}(\mathbf{S}^{-1}\mathbf{A}), u)\) is a solution of the universal mapping problem (relative to A and ST) considered in no. 1. This shows the existence and uniqueness of \(j\) .

\item
  The proof is completely analogous with that of (i), using in this case no. 2, Proposition 3, and is left to the reader.
\end{enumerate}

PROPOSITION 8. Let A be a ring and S, T two multiplicative subsets of A such that \(\mathbf{S} \subset \mathbf{T}\) . The following properties are equivalent:

\begin{enumerate}
\def\labelenumi{(\alph{enumi})}
\item
  The homomorphism \(i_{\mathbf{A}}^{\mathrm{T},\mathrm{S}}: \mathrm{S}^{-1}\mathrm{A} \to \mathrm{T}^{-1}\mathrm{A}\) is bijective.
\item
  For every A-module M, the homomorphism \(i_{\mathbf{A}}^{\mathrm{T},\mathrm{S}}: \mathrm{S}^{-1}\mathrm{M} \to \mathrm{T}^{-1}\mathrm{M}\) is bijective.
\item
  For all \(t \in \mathbf{T}\) , there exists \(a \in \mathbf{A}\) such that \(at \in \mathbf{S}\) (in other words, every element of \(\mathbf{T}\) divides an element of \(\mathbf{S}\) ).
\item
  Every prime ideal which meets T meets S.
\end{enumerate}

It has been seen above that \(i_{\mathbf{A}}^{\mathrm{T},\mathrm{S}} = 1_{\mathbb{M}} \otimes i_{\mathbf{A}}^{\mathrm{T},\mathrm{S}}\) , which immediately proves the equivalence of (a) and (b). Set \(\mathrm{T}' = i_{\mathbf{A}}^{\mathrm{S}}(\mathrm{T})\) ; then (Proposition 7) \(\mathrm{T}^{-1}\mathbf{A}\) is identified with \(\mathrm{T}'^{-1}(\mathrm{S}^{-1}\mathbf{A})\) and (a) is equivalent to saying that the elements of \(\mathrm{T}'\) are invertible in \(\mathrm{S}^{-1}\mathbf{A}\) (no. 1, Remark 5). Now, to say that \((t/1)(a/s) = 1/1\) ( \(t \in \mathrm{T}\) , \(a \in \mathrm{A}\) , \(s \in \mathrm{S}\) ) means that there exists \(s' \in \mathrm{S}\) such that \(tas' = ss'\) , which shows the equivalence of (a) and (c). We show that (d) implies (c). Let \(t\) be an element of \(\mathrm{T}\) and suppose that \(t/1\) is not invertible in \(\mathrm{S}^{-1}\mathbf{A}\) ; then there exists a maximal ideal \(m'\) of \(\mathrm{S}^{-1}\mathbf{A}\) containing \(t/1\) (Algebra, Chapter I, § 8, no. 7, Theorem 2) and \(p = (i_{\mathbf{A}}^{\mathrm{S}})^{-1}(m')\) is a prime ideal of \(\mathbf{A}\) containing \(t\) and not meeting \(\mathbf{S}\) (since the image under \(i_{\mathbf{A}}^{\mathrm{S}}\) of an element of \(\mathbf{S}\) is invertible). Conversely, if there exists a prime ideal \(p\) which meets \(\mathrm{T}\) without meeting \(\mathrm{S}\) , then no element of \(p \cap \mathrm{T}\) can divide an element of \(\mathrm{S}\) ; this proves that (c) implies (d) and completes the proof.

It follows from Proposition 8 that, amongst the multiplicative subsets T of A containing S and satisfying the equivalent conditions of Proposition 8, there exists a greatest, consisting of all the elements of A which divide an element of S (cf.~Exercise 1).

PROPOSITION 9. Let I be a right directed preordered set, \((\mathrm{S}_{\alpha})_{\alpha \in \mathbf{I}}\) an increasing family of multiplicative subsets of a ring A and \(S = \bigcup_{\alpha \in I} S_{\alpha}\) . We write \(\rho_{\beta \alpha} = i_{A}^{S_{\beta}, S_{\alpha}}\) for \(\alpha \leqslant \beta\) , \(\rho_{\alpha} = i_{A}^{S, S_{\alpha}}\) . Then \((S_{\alpha}^{-1}A, \rho_{\beta \alpha})\) is a direct system of rings and, if, for all \(\alpha \in I\) , \(\rho_{\alpha}'\) is the canonical mapping of \(\mathbf{S}_{\alpha}^{-1}\mathbf{A}\) to \(\lim_{\longrightarrow}\mathbf{S}_{\alpha}^{-1}\mathbf{A}\) , there exists a unique isomorphism \(j\) from \(\lim_{\longrightarrow}\mathbf{S}_{\alpha}^{-1}\mathbf{A}\) to \(\mathbf{S}^{-1}\mathbf{A}\) such that \(j\circ \rho_{\alpha}^{\prime} = \rho_{\alpha}\) for all \(\alpha \in \mathbf{I}\) .

For \(\alpha \leqslant \beta \leqslant \gamma\) , \(\rho_{\gamma \alpha} = \rho_{\gamma \beta} \circ \rho_{\beta \alpha}\) (no. 1, Corollary 3 to Proposition 2), hence \((S_{\alpha}^{-1}A, \rho_{\beta \alpha})\) is a direct system. We write \(A' = \varinjlim S_{\alpha}^{-1}A\) ; as \(\rho_{\alpha} = \rho_{\beta} \circ \rho_{\beta \alpha}\) for \(\alpha \leqslant \beta\) (no. 1, Corollary 3 to Proposition 2), \((\rho_{\alpha})\) is a direct system of homomorphisms and \(j = \varinjlim \rho_{\alpha}\) is the unique homomorphism from \(A'\) to \(S^{-1}A\) such that \(j \circ \rho_{\alpha}' = \rho_{\alpha}\) for all \(\alpha \in I\) . The homomorphisms \(\rho_{\alpha}' \circ i_{A^{\alpha}}^{S}: A \to A'\) are all equal, for \(\rho_{\beta \alpha} \circ i_{A^{\alpha}}^{S} = i_{A^{\beta}}^{S}\) for \(\alpha \leqslant \beta\) ; let \(u\) be their common value. Clearly the elements of \(u(S)\) are invertible in \(A'\) , which shows that there exists a homomorphism \(h: S^{-1}A \to A'\) such that \(h \circ i_{A}^{S} = u\) (no. 1, Proposition 1). Then

\[
j \circ h \circ i _ {\mathrm{S}} ^ {\mathrm{A}} = j \circ u = j \circ \rho_ {\alpha} ^ {\prime} \circ i _ {\mathrm{S} ^ {\alpha}} ^ {\mathrm{A}} = \rho_ {\alpha} \circ i _ {\mathrm{A} ^ {\alpha}} ^ {\mathrm{S}} = i _ {\mathrm{A}} ^ {\mathrm{S}}
\]

for all \(\alpha \in \mathbf{I}\) and consequently \(j \circ h\) is the identity automorphism of \(\mathbf{S}^{-1}\mathbf{A}\) . On the other hand, for all \(\alpha \in \mathbf{I}\) ,

\[
h \circ j \circ \rho_ {\alpha} ^ {\prime} \circ i _ {A ^ {\alpha}} ^ {S} = h \circ \rho_ {\alpha} \circ i _ {A ^ {\alpha}} ^ {S} = h \circ i _ {A} ^ {S} = u = \rho_ {\alpha} ^ {\prime} \circ i _ {A ^ {\alpha}} ^ {S},
\]

whence \(h \circ j \circ \rho_{\alpha}' = \rho_{\alpha}'\) for all \(\alpha \in I\) ; it follows that \(h \circ j\) is the identity automorphism of \(A'\) and consequently \(j\) is an isomorphism.

COROLLARY. Under the hypotheses of Proposition 9, let M be an A-module. We write \(f_{\beta \alpha} = i_{\mathrm{M}}^{\mathrm{S}_{\beta},\mathrm{S}_{\alpha}}\) for \(\alpha \leqslant \beta, f_{\alpha} = i_{\mathrm{M}}^{\mathrm{S},\mathrm{S}_{\alpha}}\) for all \(\alpha \in \mathbf{I}\) and let \(f_{\alpha}'\) be the canonical mapping from \(\mathrm{S}_{\alpha}^{-1}\mathrm{M}\) to \(\lim \mathrm{S}_{\alpha}^{-1}\mathrm{M}\) ; then there exists an \(\mathrm{S}^{-1}\mathrm{A}\) -isomorphism \(g\) of \(\mathrm{S}^{-1}\mathrm{M}\) onto \(\lim \mathrm{S}_{\alpha}^{-1}\mathrm{M}\) such that \(g \circ f_{\alpha} = f_{\alpha}'\) for all \(\alpha \in \mathbf{I}\) .

The corollary follows immediately from the definitions \(S_{\alpha}^{-1}M = M \otimes_{A} S_{\alpha}^{-1}A\) and \(S^{-1}M = M \otimes_{A} S^{-1}A\) and the fact that taking direct limits commutes with tensor products (Algebra, Chapter II, § 6, no. 3, Proposition 7).

